\section{Conclusions}

In this work a refocused tensorial framework has been established for the
analysis and transport of centrifugal distortion constants in rotational
spectroscopy. The central result is that the conventional quartic and
sextic centrifugal distortion sectors belong to a common linear
representation-equivalent tensor space. After pseudoinverse gauge fixing,
Watson constants can be lifted to tensor representatives, transported
linearly between principal-axis representations and Hamiltonian
reductions, and reprojected onto the desired Watson parametrization.

This common structure is particularly significant on the sextic side.
The present analysis shows that the standard sextic geometry sector
closes without introducing second derivatives of the inertia tensor as
primitive tensorial objects. Instead, the sextic geometry contributions
can be expressed entirely within the same first-level tensorial language
that governs the quartic sector, supplemented by Coriolis couplings and
the explicit cubic anharmonic contribution. In this sense, the standard
quartic and sextic centrifugal distortion sectors share the same
underlying tensorial structure and differ only in the way these
first-level objects enter the perturbative expansion.

In the formulation adopted here, the sextic equations are also made more
explicit than in their usual black-box form by separating the geometric
sector organized by $c_1$, $c_2$, and $\tau$ from the cubic-force
contribution. This opens the way to a further decomposition of the cubic
piece into semi-diagonal and genuinely three-index parts, which is of
direct practical interest for medium-sized molecules.

Within this unified baseline the first departure from the standard sector
can also be identified explicitly. On the quartic side, the channel
$H_{22}$ represents the first contribution that requires a genuinely new
intrinsic tensorial object beyond the standard first-level structure.
On the sextic side, the corresponding first post-standard candidate
appears through the term linear in $\tau(H_{22})$ in the sextic geometry
functional. These quantities are not introduced here as complete
higher-order corrections, but rather as low-cost tensorial diagnostics of
the breakdown of the standard linear transport sector.

The key point is not that $H_{22}$ must be the numerically largest
post-standard contribution. From the quantum-chemical point of view, the
dominant beyond-standard term is generally the largest one in absolute
value. From the spectroscopic fitting point of view, however, the crucial
term is the first one that breaks the linear equivalence of Watson
parameter sets across representations. Within the present framework,
fits in different representations remain physically equivalent only so
long as the $H_{22}$ diagnostic is negligible.

This diagnostic interpretation has direct practical significance. The
quantities entering $H_{22}$ and its sextic analogue are already
accessible from the same rovibrational information used in standard
calculations and therefore provide indicators of representation
sensitivity and of the need to move beyond the standard perturbative
treatment without requiring additional quantum-chemical input. The
comparison between water and formaldehyde illustrates this behaviour:
the quartic $H_{22}$ diagnostic is significantly larger in the case
where VPT-based constants show stronger deviations from higher-level
results.

The present work therefore establishes the appropriate baseline for
subsequent developments. It clarifies the common tensorial structure of
standard quartic and sextic centrifugal distortion, makes explicit the
role of pseudoinverse gauge fixing in the inverse problem, and
identifies the first quartic and sextic diagnostics of the breakdown of
the standard transport picture. A more complete treatment of the
post-standard channels and their perturbative organization will be the
subject of future work.

In this sense, the present framework should not be viewed as a mere
reparametrization of Watson constants. Its central achievement is the
identification of the minimal linear transport sector underlying the
standard quartic and sextic constants together with the first
contribution that forces one to leave that sector.

The main contributions of the present work may therefore be summarized
succinctly as follows. First, standard quartic and standard sextic
centrifugal distortion are shown to belong to a common linear tensor
sector that can be transported between representations. Second, the
validity of this linear sector can be tested at essentially harmonic
cost through the quartic $H_{22}$ diagnostic and its sextic analogue.
Third, when the genuine three-index cubic sector is small, the full set
of distortion constants becomes accessible at roughly twice the cost of
a harmonic calculation based on analytic gradients. Fourth, the
distortion constants are reformulated in a way that separates geometric,
harmonic, and cubic contributions, allowing them to be treated at
different computational levels. Fifth, the whole formalism is made
operational through an app and a Python implementation, so that tensor
transport and first-breakdown diagnostics can be used directly in
practical spectroscopic workflows. This last point is especially
relevant on the sextic side, where standard constants are presently
available in only a very limited number of quantum-chemical
implementations and no widely used representation-transform machinery
exists.

To facilitate practical use of the framework, the corresponding
transformations and diagnostics have been implemented in a dedicated
computational tool available either as a macOS application or as a
Python code. The analyses discussed in the present work can be
reproduced through the app workflow, which therefore serves both as a
dissemination tool and as an operative implementation of the
spectroscopic protocol introduced here. The program implements the
standard quartic and sextic transport machinery together with the first
post-standard diagnostics, namely the quartic $H_{22}$ term and its
sextic linear analogue, when harmonic input is provided (`.fchk` or
`xyz + Hessian`).

The tensor transport perspective proposed here provides a transparent
interpretation of the representation dependence of centrifugal
distortion constants and places quartic and standard sextic terms
within a common structural framework. Beyond its conceptual relevance,
this formulation offers a practical route for reconstructing and
validating distortion constants once the underlying tensor
representatives are known.

The tensor formulation developed in this work also provides a
structural interpretation of centrifugal distortion constants and of
their representation dependence within Watson's effective rotational
Hamiltonian. In this framework, quartic and standard sextic
centrifugal distortion constants emerge as elements of a common linear
tensor sector generated by the first derivative of the inverse inertia
tensor and by Coriolis couplings. Once the gauge of the inverse
problem is fixed through Moore--Penrose pseudoinversion,
transformations between Hamiltonian representations reduce to linear
transport operations in this tensor space.

The analysis also clarifies the origin of the first departure from
this linear structure. While the linear tensor sector accounts for
both quartic and standard sextic distortion constants, the appearance
of cubic force constants introduces genuinely new contributions and
therefore marks the first nonlinear layer of the expansion. In this
sense, the tensor formulation not only provides a compact
interpretation of representation transformations, but also exposes the
geometric structure underlying centrifugal distortion constants and
reveals a natural hierarchical organization of centrifugal distortion
terms.

From a spectroscopic perspective, the present tensor formulation also
provides a natural interpretation of the transformations between
different Hamiltonian representations, such as the well-known $A$ and
$S$ reductions. Within the linear tensor sector identified here, these
transformations correspond to linear transport operations acting on the
underlying tensor representatives, followed by projection onto the
parameter space of the chosen Hamiltonian form.

From this perspective, centrifugal distortion constants should no
longer be viewed as independent parameters attached to a specific
Hamiltonian representation, but rather as representation-dependent
projections of a smaller set of underlying tensor quantities.

The present tensor formulation therefore provides a structural
interpretation of quartic and standard sextic centrifugal distortion
constants as elements of a common linear tensor sector. Beyond its
conceptual implications, this viewpoint may simplify the analysis of
representation transformations and could potentially facilitate the
reconstruction of distortion constants from reduced sets of geometric
derivatives. Together with the recent developments of our group in
obtaining accurate equilibrium structures at DFT cost and vibrational
corrections at essentially harmonic cost, this framework also makes it
realistic to guide and interpret spectroscopic fits for molecules in the
30--50 atom range with errors already comparable to those achieved by
highly accurate methods on much smaller systems.

The same viewpoint also clarifies that, at the harmonic and
standard-VPT2 levels, the Cartesian-normal-mode and curvilinear
internal-coordinate formalisms are not competing descriptions but
coordinate realizations of the same underlying rovibrational tensor
structure.

More broadly, the gauge-fixed transport picture developed here suggests
that the effective rovibrational Hamiltonian itself may be viewed as a
quotient structure of tensorial representatives modulo contact-
transformation freedom, with quartic and sextic centrifugal distortion
providing a first explicit and computationally accessible realization
of this more general viewpoint.

From a broader mathematical perspective, the inverse reconstruction
problem treated here is also closely related to a structural
identifiability problem: Watson constants do not uniquely determine an
underlying tensor representative, but only an equivalence class, and
the Moore--Penrose reconstruction provides the canonical minimum-norm
choice within that class.
